\chapter{Deques}\label{chap:deques}

A deque (double-ended queue) is a linear data structure that supports insertion and removal at both ends in $O(1)$ time. It serves as a fundamental building block for other linear data structures like lists, stacks, and queues.

\section{Design rationale}

We implemented the deque as a doubly-linked circular list with a sentinel node. The sentinel node simplifies the implementation by eliminating edge cases (such as empty lists) and avoiding the need for null checks during insertion and removal operations.

The deque implementation uses intrusive data structures for genericity. This means that the list links (`next` and `prev` pointers) are embedded directly within the user's data structure. This approach avoids dynamic memory allocation for list nodes and allows an element to belong to multiple lists simultaneously (by having multiple link fields).

The deque header contains the sentinel node and metadata about the deque, such as its size.

\section{Deque interface}\label{deques:interface}

The interface provides macros and functions for initializing deques, pushing/popping elements at both ends, and iterating over the elements.

\begin{code}
  \lstinputlisting[style=C,linerange={6-36}]{libellul/include/libellul/structure/deque/deque-export-def.h}
  \caption{The generic deque interface.}
  \label{code:deque:interface}
\end{code}

\section{Benchmark}

We benchmarked the performance of our deque implementation (linked list based) against a dynamic array implementation.

\begin{figure}[H]
    \centering
    \subfloat[\texttt{push\_front}]{\includegraphics[width=0.45\textwidth]{deque_push_front.png}}
    \hfill
    \subfloat[\texttt{push\_back}]{\includegraphics[width=0.45\textwidth]{deque_push_back.png}}
    \\
    \subfloat[\texttt{pop\_front}]{\includegraphics[width=0.45\textwidth]{deque_pop_front.png}}
    \hfill
    \subfloat[\texttt{pop\_back}]{\includegraphics[width=0.45\textwidth]{deque_pop_back.png}}
    \caption{Performance comparison of Linked Deque vs. Array Deque.}
    \label{fig:deque:bench}
\end{figure}
